
The corpus protocol is \texttt{PROTOCOL.md} and the practice pilot's is
\texttt{AUDIT-PROTOCOL-2.md}; both are in the archive and both were committed
before the first result they govern existed. This appendix lists what changed
afterwards.

\paragraph{Why amendments are reported rather than folded in}
A pre-registration that is quietly edited is worse than none, because it
claims a discipline it does not have. Each amendment below names its date, the
defect that forced it, and the result it moved.

\begin{description}
\item[Amendment 1 --- a register's values must recur across time.]
Registered rule 3.3 admits any attribute reused across two cases on average.
On BPI Challenge 2019 that admits the purchasing-document number, which is a
document identifier and not a maintained entity: almost none of its values
recur across the temporal split. The amendment adds a second, outcome-free
condition --- at least \minTrainCoverage\ of test cases must carry a value
seen in training --- and \emph{both} selections are reported for every log.

\item[Amendment 2 --- ``any timestamp'' had to be implemented.]
Rule 3.1 excludes timestamps. Matching a fixed list of names did not implement
that: it missed five differently-spelled date columns, two of which the
registered rules then made the high-cost entity. A column is now a timestamp
if its name looks like one or if more than $90\%$ of its non-missing values
parse as a date. Both tests read feature values only.

\item[Amendment 3 --- a relabelling of the activity is the activity.]
Some logs carry a second column that is a one-to-one recoding of
\texttt{concept:name}. Admitting it as a feature admits the activity sequence,
which is not available at case creation.

\item[Amendment 4 --- coarsenings of the register get their own role.]
Type, subtype and service are deterministic functions of the item. Sweeping
them into the intake block makes the intake block partly a register, and the
increment then measures the residual rather than the field. They are a
separate role and \ref{app:layers} uses it.

\item[Amendment 5 --- the caused-by family is an outcome.]
A caused-by attribution is the output of diagnosis, not an intake observation.
Without this the registered rules made the caused-by configuration item the
high-cost entity on the flagship log.

\item[Amendment 6 --- the pilot frame's topical filter was nulled.]
The frame's keyword filter was applied to titles and abstracts, which is a
filter on how a paper describes itself rather than on what it does.
\texttt{results/s06\_frame.csv} reports what nulling it does to the estimates.
\end{description}

\paragraph{What the protocol fixed, before any of the results it governs}
The design space of Section~\ref{sec:surface}; the reference specification;
the exclusion of the intercept-only rung from every admissible set; the six
register-quality mechanisms and their levels; the five planned contrasts and
the Holm procedure; the two families over which simultaneous bands are
taken; and the two-stratum design of the practice pilot, including the
decision to adjudicate a sample of the papers the screen \emph{rejected}. The
commit that fixed them adds \texttt{scripts/spec.py} and
\texttt{PLAN-ROUND19.md} and no result file.

\subsection{Related work, strand by strand}
\label{app:related}

Four further strands bear on the argument. \emph{Many-analysts studies} put
one question and one dataset to independent teams and observe the spread
directly \citep{silberzahn2018many,breznau2022observing}: they establish that
analytic variability is large without supplying an estimand for it, which is
what the surface is. \emph{Predictive multiplicity} and Rashomon sets study
near-equally accurate models that disagree on individuals
\citep{marx2020predictive,semenova2022existence} --- multiplicity at the level
of a fitted model, where this paper's is at the level of the design that
produced it. \emph{Benchmark-variance work in predictive process monitoring}
reports how far results move with the pipeline and the protocol
\citep{ramamaneiro2021deep}, which is the evidence this paper's decomposition
partitions. And the \emph{sensitivity-analysis} tradition of bounding an
unmeasured assumption \citep{vanderweele2017sensitivity} is the form
Section~\ref{sec:quality}'s leakage tipping point takes. On the estate side,
the accuracy of a configuration management database is itself a studied
quantity \citep{suriadi2017imperfection,mohammed2025dataquality}, and the
layered service models practitioners maintain are why
Section~\ref{sec:programme} asks which \emph{layer} of a register pays.
Supplement~\ref{app:related} sets out each strand and what this paper takes
from it.


%%  ROUND TWENTY-FIVE.  Moved from Section 2 for length, unchanged.

\paragraph{Variable importance and incremental value}
That a feature's apparent importance depends on what else is in the model is
the standard motivation for conditional and permutation-based importance
\citep{strobl2008conditional,fisher2019models}, for Shapley-value attributions
\citep{lundberg2017unified,covert2020sage}, and for nonparametric formulations
of importance as a population quantity \citep{williamson2021vimp}, which
average over subsets precisely because a single conditioning set is arbitrary.
The clinical-prediction literature has made the same point about the
\emph{incremental} value of a marker for two decades: the increase in the area
under the ROC curve from adding a predictor depends on how good the existing
model is \citep{pencina2008evaluating,cook2007roc}, is a poor guide to
clinical usefulness \citep{cook2007roc,vickers2008decision}, and is routinely
misread \citep{kerr2011evaluating}. What that literature establishes, and what
we take as given, is that the baseline and the metric are choices. What it
does not supply is an object that reports the whole dependence, an estimator
for it, or inference over it.

\paragraph{Decision-analytic evaluation}
Decision curve analysis \citep{vickers2006dca,vickers2016net} makes the
operating point explicit by reading net benefit across an exchange rate, and
is the reason the operating point is an axis here rather than a footnote.
Guidance on reading such a curve is well developed \citep{vickers2019simple};
what is less settled is what a claim about a \emph{band} of thresholds
requires, which is why the bands in Section~\ref{sec:dca} are simultaneous
rather than pointwise.

\paragraph{Specification sensitivity}
That analytic conclusions move with defensible analytic choices is the subject
of multiverse analysis \citep{steegen2016increasing}, specification-curve
analysis \citep{simonsohn2020specification}, vibration of effects
\citep{patel2015assessment} and, in machine learning evaluation, of work on
underspecification \citep{damour2022underspecification} and on the fragility
of benchmark rankings \citep{bouthillier2021variance}. Our design space and
our figures are recognisably in that tradition. What this paper adds to it is
a decomposition of the surface's variance into axis contributions,
simultaneous inference over the surface rather than a distribution of point
estimates, region labels with an explicit decision rule, and a regret
benchmark against the conventional report. Specification-curve analysis, as
usually practised, displays the curve and applies a permutation test to its
median; it does not tell a reader which sub-region of the design space
supports a sign.

\paragraph{Data valuation}
Shapley-based data valuation \citep{ghorbani2019data,jia2019towards} values
\emph{instances} against a fixed learner and metric; the axes we vary are
exactly the ones held fixed there. The two are complementary: a data value is
itself a point on a specification surface.

\paragraph{Which axis a study varies}
Some studies do vary one axis deliberately: the metric axis in effort-aware
defect prediction \citep{carka2022effortaware}, the encoding axis in
remaining-time prediction \citep{liu2025encoding}, the inter-case feature
block in predictive process monitoring \citep{senderovich2017intercase}. What
is missing is an object that holds several axes at once and says what their
joint variation implies for a sign.

\paragraph{Prediction on incident event logs}
The empirical setting is predictive process monitoring, surveyed and
benchmarked across real logs
\citep{teinemaa2019outcome,verenich2019survey} on data of the kind process
mining formalises \citep{vanderaalst2016process}. We predict at case creation
from creation-time attributes with no prefix, which is tabular classification
on a process log rather than prefix-based monitoring, and that difference
bounds what the corpus result transfers to. The benchmark protocols in that
literature fix a learner, an encoding and a metric per experiment, and this
paper measures the sensitivity of a conclusion to \emph{those} fixings. It
does not measure sensitivity to \textbf{prefix length, prefix bucketing or
sequence encoding}, which are the axes that literature varies most, because
none of them exists at a single prediction point; Section~\ref{sec:lim}
carries the transfer as an assumption rather than a result. Configuration
management databases and their data quality are treated in the IT service
management literature as an operational problem \citep{marrone2011itil}; we
are not aware of an evaluation of one against a named prediction task with the
design space held open.

\paragraph{Data quality as an axis}
The effect of data quality on model performance has been studied directly on
tabular data \citep{mohammed2025dataquality} and, in process mining, for
low-quality activity labels \citep{comuzzi2025labels}. Those studies vary
quality with the feature set fixed; we vary quality \emph{as one axis of the
same surface} on which the feature set and the encoding are others, which is
what makes the interaction between them visible ---
Section~\ref{sec:quality} reports a quality defect that is exactly invisible
under one encoder and not under another. Leakage, which is availability
misspecified, has its own literature
\citep{kaufman2012leakage,kapoor2023leakage,weytjens2022leakage}, and the
decision time of Section~\ref{sec:twodecision} makes it a declared choice
rather than a mistake.

\subsection{The measure on the axes, and the two axes that are not analysis choices}
\label{app:measurechoice}

%%  ROUND TWENTY-FIVE.  Moved from the article for length, unchanged.

\paragraph{On the measure}
A design space is a set; a decomposition, a weighted mean or a robustness
index needs a \emph{distribution} on it, and ``uniform over the cells I
enumerated'' is not neutral --- it moves if a level is duplicated, if five
folds are written as five levels rather than one, or if a grid is refined.
This paper separates the weight on an axis from the weights within it,
declares three product measures and an envelope over a fourth
(Section~\ref{sec:sobol}), and reports every weighted quantity under all of
them.

Two of these axes deserve comment because they are usually not treated as
design choices at all.

The \textbf{decision time} $\tau$ fixes which fields exist when the prediction
is made. It is not the same as the split rule and it is not a data-cleaning
detail: moving $\tau$ later in a process admits fields written in between, and
Section~\ref{sec:cmdb} is a case in which that single move takes a field's
increment from resolvably positive to indistinguishable from zero.

The \textbf{register quality} $q$ acts on $f$ rather than on the comparison. A
field that is a lookup into a maintained register --- a configuration item, a
customer master record, a product catalogue --- inherits that register's
completeness and accuracy. We implement \nQualityKinds\ mechanisms besides the
clean field, each parameterised by the training half alone: population loss
with the long tail absent, with the core absent, and at random; identity
error, in which a share of rows carry a wrong value drawn from the training
alphabet; reconciliation failure, in which a share of identities exist twice;
and staleness, in which an identity is present only if the register had
already seen it by the split point. The first three model coverage, the next
two accuracy, the last discovery; the taxonomy follows the event-log
imperfection patterns of \citet{suriadi2017imperfection} where they apply.

Real registers do not fail at random, so three \emph{further} mechanisms make
the missingness depend on something: on \emph{time}, where a row's probability
of carrying an identity rises with its position in the log; on \emph{content},
where coverage depends on the item's recorded type through a map estimated on
the training half; and on a \emph{propensity} built from the intake block,
which makes missingness correlated with the outcome without ever reading a
test outcome. Every mechanism is swept over \nQualityLevels\ severities and
every stochastic one repeated over \nQualitySeeds\ seeds, and each is verified
by execution to be a function of the training half alone --- perturbing the
test half and re-degrading must leave the training half's degraded values
identical, which \texttt{results/s01\_trainonly.csv} records over \nTrainOnly\
checks.

\subsection{Targets, splits and the exclusion codes}
\label{app:targets}

%%  ROUND TWENTY-FIVE.  Moved from the article for length, unchanged.

\subsection{Targets, splits and exclusions}

Two targets are defined for every log: \emph{handover}, whether more than one
resource group touches the case, and \emph{duration}, whether the case runs
longer than the training half's median. Cases are ordered by their first
timestamp and split $70/30$; the rolling-origin scheme adds \nFolds\
expanding-window folds. Fields that are determined only at closure are excluded by name, which is the
standard leakage discipline \citep{kaufman2012leakage,kapoor2023leakage} and,
for event logs specifically, the discipline of
\citet{weytjens2022leakage}. A log--target pair is excluded only for a
registered reason --- too few cases, no variation in the test half, a prevalence outside
$[\prevLo, \prevHi]$, or the absence of one of the three roles --- and every
exclusion is reported with its code in Table~\ref{tab:corpus}, which accounts
for every downloaded log that is not among the admitted pairs --- including
the held-out set, \nHeldOutLogs\ of the checksummed files, downloaded for a
separate test and never offered to these rules at all.


\textbf{Every admitted pair is analysed and reported.} Running the full
surface only on the logs whose \emph{reduction} is resolvable would condition
the analysis on the outcome. The surface is computed on all \nPairs\ admitted
pairs through one code path, and the pairs on which the register is worth
nothing are in the same tables as the rest.

\subsection{The admission margin at the upper edge}
\label{app:margin}

%%  ROUND TWENTY-FIVE.  Moved from the article for length, unchanged.

\paragraph{The pair carrying the most cells sits closest to the admission
edge} \logAdmittedMaxPrev\ has the highest prevalence of any admitted pair,
\prevAdmittedMax, against the window's upper bound of \prevHi; the nearest
pair the headroom rule turns away, \logExcludedAboveMin, is at
\prevExcludedAboveMin, and \nExcludedAbove\ pairs in all are excluded above
the bound. That pair is also the largest single contributor to the corpus,
carrying \nCellsLargestPair\ of the \nComputationalCells\ admissible scalar
cells. A rule with a sharp edge decides some case narrowly, and this is that
case: the corpus's biggest component is admitted on a margin narrower than
the gap to the pairs excluded beside it. Table~\ref{tab:family} reports every
headline statistic on the ITSM-like subset as well, which is the population
that does not turn on this decision.

